%! Justifying a Claim About the Difference of Two Means Based on a CI — Unit 7, Lesson 7.7 — Exit Ticket
\documentclass[10pt]{article}
\usepackage{apstats-article}
\usepackage{apstats-boxes}
\newcommand{\TallMath}[1]{$\displaystyle #1\rule[-1.4em]{0pt}{3.2em}$}

\begin{document}

\pageheader{Unit 7, Lesson 7.7}{Exit Ticket}
\namedateperiod

\vspace{0.15in}

A researcher compares mean resting heart rate (bpm) for runners and non-runners.
Using independent random samples of 40 runners and 40 non-runners,
a 95\% confidence interval for $\mu_{\text{runner}} - \mu_{\text{non-runner}}$ is
$(-18.4,\; -6.2)$ bpm. All conditions for inference were verified.

\begin{enumerate}
  \item Interpret the confidence interval in context.
        \writelines{2}

  \item A cardiologist claims that runners have a lower mean resting heart rate than
        non-runners. Does the confidence interval provide convincing evidence for this
        claim? Explain by referencing the interval.
        \writelines{2}

  \item If the researcher had used samples of 160 runners and 160 non-runners
        (all other values the same), how would the width of the CI change?
        Circle one: \quad \textbf{Narrower} \quad / \quad \textbf{Wider} \quad / \quad \textbf{Same} \\[4pt]
        Explain: \writeline
\end{enumerate}

\end{document}
